% Limitations section

\paragraph{Single model, synthetic tasks.}
All experiments use MiniMax M2.5 on a synthetic task corpus (300 tasks, five themes, three difficulty levels).
The output-driven efficiency mechanism may be sensitive to model-specific instruction-following behavior, and the 100\% task success rate means we measure token efficiency and retrieval quality but not response quality differences.
Validating these findings on frontier models and real user queries remains necessary.

\paragraph{Limited exploration budget.}
Thompson Sampling with 12 arms and 300 episodes per condition yields roughly 25 episodes per arm under uniform exploration.
C2 and C3 reach a stable selection pattern by approximately task 33--48 (changepoint analysis), with \texttt{pure\_relevance} dominating thereafter; however, the posterior continues to sharpen through episode 300.
C4's compressed posteriors (max $E[\theta]{=}0.472$) may reflect insufficient data rather than a fundamental property of anchor-based rewards; longer runs would clarify this.

\paragraph{No human quality evaluation.}
Our reward hacking assessment relies on rating stability (no detectable drift across 1{,}200 episodes) rather than independent quality judgment.
We cannot confirm that efficiency gains preserve output quality; the stability finding is consistent with reward hacking avoidance but does not prove it.
More broadly, token efficiency is a proxy metric: shorter outputs are not necessarily better outputs, and without human evaluation of response quality, the 12--20\% token reduction could reflect less thorough responses rather than more concise ones.

\paragraph{Confounded cross-experiment comparison.}
Experiments~1 and~2 differ in multiple design dimensions simultaneously (library size, skill specificity, number of presets, task difficulty distribution), so improvements cannot be attributed to a single factor.
The Jaccard overlap reduction from ${>}$0.90 to 0.52--0.59 confirms retrieval differentiation, but the overlap remains moderate.
Additionally, C4's anchor embedding parser compresses the recency dimension (mean 0.136 vs.\ importance 0.62 and relevance 0.55), effectively reducing feedback to two dimensions and limiting the bandit's ability to distinguish recency-weighted presets.

\paragraph{Static skill library and parser specificity.}
Skills were fixed throughout each experiment; a production system with evolving skills would face non-stationarity that our bandit formulation does not address.
The feedback parser is also specific to MiniMax's output patterns (parse rates: 93--97\%), and the 5--7\% parse failure rate introduces noise into the reward signal.

\paragraph{Single experimental run.}
Results are from a single execution per condition with no repeated runs under different random seeds.
We cannot estimate seed-dependent variance or confirm reproducibility of the specific convergence patterns (e.g., C3's lock-in to \texttt{pure\_relevance} by task 44).
The fixed-seed task shuffling ensures deterministic condition assignment but does not address the inherent stochasticity of LLM generation and Thompson Sampling exploration.

\paragraph{Single-annotator ground truth.}
Ground truth skill annotations (1--3 relevant skills per task) were created by a single annotator.
Inter-annotator agreement was not assessed, and the subjective nature of skill relevance means that retrieval quality metrics (NDCG@5, MRR, hit rate) reflect one person's judgments rather than a consensus standard.

\paragraph{Missing ablation controls.}
The experimental design lacks a rubric-only ablation (appending the feedback rubric without any bandit updates) and a comparison against standard MAB baselines (e.g., UCB1, $\epsilon$-greedy).
Without a rubric-only control that explicitly disables the bandit, the rubric prompt effect and adaptive retrieval effect cannot be fully disentangled.
Additionally, C3's explanation embedding storage constitutes a potential confound: C3 maintains additional semantic information not available to C2, so its superior retrieval quality may partly reflect richer stored context rather than better reward signals alone.

\paragraph{Dimension label semantics.}
The ``importance'' dimension is computed as a Laplace-smoothed historical success rate, which more closely tracks usage frequency than intrinsic importance.
This labeling may affect how the LLM interprets the feedback rubric, potentially conflating frequency-based and quality-based assessments in its dimension ratings.

\paragraph{Narrow evaluation domain.}
All tasks and skills are drawn from a single domain (AI/ML research), limiting generalizability to other knowledge domains where the relative informativeness of recency, importance, and relevance dimensions may differ substantially.

\paragraph{Narrow architectural evaluation.}
All experiments operate within the vector-based RAG paradigm.
Our 205-skill library is small enough that agentic keyword search, long-context models, or fine-tuning could plausibly match or exceed our retrieval quality without the bandit machinery.
We do not evaluate whether the Thompson Sampling mechanism or the self-assessment reward signal transfers to these alternative architectures, though the mechanisms are not structurally dependent on vector retrieval (see Section~\ref{sec:discussion}).

\paragraph{Related work scope.}
The related work section does not discuss connections to learning-to-rank (LTR) methods, which optimize ranking functions via supervised or semi-supervised learning, or to the metacognitive monitoring literature in cognitive science, which provides a theoretical grounding for the input-assessment concept.
These connections represent directions for future positioning.
